% results.tex (v14, 2026-10-20): tightened wording, concise captions (details in Methods), Fig. 6 paragraph for the v14 design. Delta = RMSE(first) - RMSE(second) on the same site-years; 95% paired block-bootstrap interval in brackets.
\section{Results}\label{sec:results}

\subsection{Effect of the year-specific climate input}
The observations are unevenly distributed in time (Fig.~\ref{fig:data}b, Table~\ref{tab:data}): \ntopyrshareAk\% of the Alaska site-years were observed in \ntopyrAk, whereas the Lena Delta labels are spread over \nnyrLe{} years. Most measurements were made in August (median day of year \ndoymedAk{} in Alaska, \ndoymedLe{} in the Lena Delta and \ndoymedCa{} in Canada), before the end of the thaw season, so that the TDD accumulated to the observation day amounted to a median of \ntddfracAk, \ntddfracLe{} and \ntddfracCa{} of the full-season sum (equation~(\ref{eq:tddobs})).

Replacing the mean climate input by the year-specific input reduced the spatial-block RMSE of the recalibrated Stefan model in the Lena Delta ($\Delta$RMSE \ndforcLePo{} cm) and that of the physics-residual model in the Lena Delta (\ndforcLeRo{} cm) and Canada (\ndforcCaRo{} cm) (Fig.~\ref{fig:accuracy}a). In Alaska the change was not resolved (Stefan \ndforcAkPo{} cm), consistent with the concentration of the labels in one year, for which the year-specific input adds little beyond the spatial pattern. Beyond accuracy, the year-specific input ties each prediction to the climate of a stated year, which the yearly maps presented below require.

\paperfig{Fig3_accuracy}{1.0}{\textbf{Climate input, validation design and conditions of the physics-residual gain.} \textbf{a}, Change in spatial-block RMSE when the 2011--2024 mean climate input is replaced by the year-specific input. \textbf{b}, RMSE of the four methods under five validation designs (upper; a value beyond the axis is printed) and the difference between the physics-residual model and the recalibrated Stefan model under each design (lower). \textbf{c}, The same difference for six covariate sets under spatial-block validation. \textbf{d}, The same difference within regional terciles of eight covariates; outlined cells have intervals excluding zero, and hatched cells have fewer than 30 site-years or five blocks. Intervals are 95\% paired block-bootstrap intervals; negative values favour the first-named model.}{fig:accuracy}

\subsection{Accuracy under spatial, temporal and regional validation}
Model rankings depended on the validation design (Fig.~\ref{fig:accuracy}b, Supplementary Table~S5). Random cross-validation gave the lowest errors to direct ML (\nrmseAkrandDz{} cm in Alaska), and spatial-block cross-validation raised the direct-ML error by \noptimAkDz, \noptimLeDz{} and \noptimCaDz{} cm in Alaska, the Lena Delta and Canada. Year hold-out raised the direct-ML error in the Lena Delta from \nrmseLeblkDz{} to \nrmseLeyrDz{} cm, and the block-by-year design produced the largest within-region errors of all methods fitted to the labels of a region. When an entire region was withheld, the direct-ML error in the Lena Delta rose to \nrmseLeregDz{} cm, against \nrmseLeregPz{} cm for the Stefan model with the coefficient of the other regions ($\Delta$RMSE \ndrefLeregDz{} cm); the physics-residual model stayed close to the Stefan model there (\ndrefLeregRo{} cm) and improved on it in Canada (\ndrefCaregRo{} cm). Spatial-block cross-validation is therefore the main design and region hold-out the test of transferability.

\subsection{Gain of physics-residual learning and its controls}
Under spatial-block cross-validation with the year-specific input, the physics-residual model reduced the RMSE of the recalibrated Stefan model from \nrmseLeblkPo{} to \nrmseLeblkRo{} cm in the Lena Delta ($\Delta$RMSE \ndrefLeblkRo{} cm) and from \nrmseCablkPo{} to \nrmseCablkRo{} cm in Canada (\ndrefCablkRo{} cm), whereas in Alaska the two models reached \nrmseAkblkPo{} and \nrmseAkblkRo{} cm (\ndrefAkblkRo{} cm). The gain depended on the covariates available to the residual learner (Fig.~\ref{fig:accuracy}c): with climate covariates only, the residual model did not improve on the Stefan model in Alaska (\ndfeatAkFclim{} cm) or Canada (\ndfeatCaFclim{} cm), and adding terrain and soil covariates produced the gains above. MODIS vegetation and snow timing or high-resolution soil and land-cover classes changed the result by less than 0.5 cm in the Lena Delta and Canada, and the latter added a small gain in Alaska (\ndfeatAkFhires{} cm).

The result also depended on how the physical model entered the learner (Supplementary Fig.~S2). Stefan pseudo-labels reproduced the Stefan model (Lena Delta \ndrefLeblkDo{} cm), and the Stefan prediction as an input feature gave results similar to direct ML. With all labels of a region, direct ML reached lower point errors than the physics-residual model (Alaska \nrmseAkblkDz{} versus \nrmseAkblkRo, Lena Delta \nrmseLeblkDz{} versus \nrmseLeblkRo{} and Canada \nrmseCablkDz{} versus \nrmseCablkRo{} cm), but its difference from the recalibrated Stefan model was resolved only in Canada (\ndrefCablkDz{} cm), and its error increased under region hold-out. Replacing CatBoost by random forest, a multilayer perceptron, TabM or FT-Transformer did not change the sign of the physics-residual gain in the Lena Delta and Canada.

Across regional terciles of eight covariates, the gain occurred in most conditions of the Lena Delta (\ngcLesig{} of \ngcLecells{} terciles with intervals excluding zero) and Canada (\ngcCasig{} of \ngcCacells) and increased with the error of the Stefan model (correlation across terciles \ngcLecorr{} in the Lena Delta and \ngcCacorr{} in Canada; Fig.~\ref{fig:accuracy}d). In the warmest TDD tercile of the Lena Delta the gain reached \ndtddLewarm{} cm, whereas no tercile in Alaska changed by more than \ngcAkmaxabs{} cm. Terrain variables accounted for \nshapLeResTerrain\% of the mean absolute SHAP value of the residual model in the Lena Delta and soil variables for \nshapCaResSoil\% in Canada (Supplementary Fig.~S3).

\subsection{Sample efficiency and transfer to new regions}
With 25 training labels drawn from the site-years of a region, the physics-residual model was at least as accurate as a Stefan model recalibrated on the same labels (Alaska \nllAkRotwfive, Lena Delta \nllLeRotwfive{} and Canada \nllCaRotwfive{} cm), and with 100 labels its gain was resolved in all three regions (\nllAkRohundred, \nllLeRohundred{} and \nllCaRohundred{} cm; Fig.~\ref{fig:labels}a). Direct ML trained on 25 labels had a higher point error than the Stefan model in Alaska (\nllAkDztwfive{} cm) and Canada (\nllCaDztwfive{} cm). Physics-input ML and physics pseudo-labels followed the patterns of their full-data results (Supplementary Fig.~S4).

\paperfig{Fig4_labels}{1.0}{\textbf{Accuracy with few labels, in new regions and in external CALM regions.} \textbf{a}, Difference in RMSE between each method and a Stefan model recalibrated on the same $n$ training labels within each region (mean of three draws; band, 95\% block-bootstrap interval of the physics-residual model). \textbf{b}, RMSE in each region for models trained on the other two regions plus $n$ labels of the target region; a value beyond the axis is printed. \textbf{c}, External CALM sites numbered by sub-region, and the label cells of the study regions. \textbf{d}, RMSE in the external regions against the number of labels from the same region, on the \nftrows{} site-years with values of all three products; horizontal lines mark the products. \textbf{e}, RMSE of the physics-residual model with no and with five local labels, and of ESA CCI v5, by sub-region (site-years in brackets).}{fig:labels}

In the transfer experiment, each region was predicted by models trained on the other two regions and $n$ labels of the target region (Fig.~\ref{fig:labels}b). Without target labels, direct ML reached \ntlrmseLeDszero{} cm in the Lena Delta, against \ntlrmseLePzzero{} cm for the source-coefficient Stefan model and \ntlrmseLeRszero{} cm for the physics-residual model; 25 target labels reduced the physics-residual error to \ntlrmseLeRstwfive{} cm and the direct-ML error to \ntlrmseLeDstwfive{} cm. In Alaska the physics-residual model gave the lowest error at every $n$ (\ntlrmseAkRszero{} cm without target labels). In Canada, direct ML trained on the source regions was the most accurate (\ntlrmseCaDszero{} cm); recalibrating the Stefan coefficient on 25 Canadian labels increased its error from \ntlrmseCaPzzero{} to \ntlrmseCaPntwfive{} cm, and the physics-residual model, which uses this anchor, followed (\ntlrmseCaRszero{} to \ntlrmseCaRstwfive{} cm), reflecting the wide range of thaw depths among the Canadian sites.

The same procedure was applied to eight external CALM regions with \nftScoredRows{} site-years not used in model development (Fig.~\ref{fig:labels}c--e, Supplementary Fig.~S5). On the \nftrows{} site-years of seven regions with values of Wei et al., Liu et al. and ESA CCI v5, the physics-residual model trained on the three study regions reached \nftAllRszero{} cm without local labels; five labels from the same region reduced the error to \nftAllRsfive{} cm and 25 labels to \nftAllRstwfive{} cm, whereas direct ML reached \nftAllDstwfive{} cm with 25 labels. With five labels the physics-residual model was more accurate than ESA CCI v5, which is not trained on ALT observations (\nftAllCcifivezero{} cm on the same site-years; $\Delta$RMSE \nftdRsCcifivefive{} cm on all \nftdCciRows{} site-years with a value). Wei et al. and Liu et al., which include these sites in their training data, reached \nftAllWeizero{} and \nftAllLiuzero{} cm. Most of the reduction came from recalibrating the Stefan coefficient; the residual learner did not reduce the error further at these sites (\nftdRsPnfive{} cm with five labels), where the active layers of West and Central Siberia and Svalbard exceed the range of the study regions. The pooled curve combines different regional responses (Fig.~\ref{fig:labels}e): with five labels the error decreased in West Siberia, Central Siberia, European Russia and Svalbard, but increased in Chukotka (\nftChRszero{} to \nftChRsfive{} cm), Kamchatka (\nftKaRszero{} to \nftKaRsfive{} cm; 16 site-years, not shown) and north-eastern Siberia (\nftNsRszero{} to \nftNsRsfive{} cm), where the error without local labels was already small.

\subsection{Comparison with published products in the same cells and years}
On the common site-years (\nyrscc), the physics-residual model reached \nccLeRoall{} cm in the Lena Delta, against \nccLeCcifiveall{} cm for ESA CCI v5, \nccLeAaltoall{} cm for Aalto et al., \nccLeCcifourall{} cm for ESA CCI v4, \nccLeLiuall{} cm for Liu et al. and \nccLeWeiall{} cm for Wei et al. (Fig.~\ref{fig:products}a, Supplementary Table~S6). On the Lena Delta site-years more than 25 km from the training sites of Wei et al. (\nnprodLeweimask{} site-years with a value of Wei et al., \nnprodLeccifivemask{} with a value of ESA CCI v5), the paired differences were \ndprodLeweiRomask{} cm relative to Wei et al. and \ndprodLeccifiveRomask{} cm relative to ESA CCI v5. In Alaska (\nccAkRoall{} versus \nccAkWeiall{} and \nccAkLiuall{} cm for Wei et al. and Liu et al.) and Canada (\nccCaRoall{} versus \nccCaWeiall{} and \nccCaLiuall{} cm), the model and the two ML products performed on par.

On the Lena Delta site-years, the two ML products assign most cells values between about 40 and 70 cm regardless of the observed ALT (Fig.~\ref{fig:products}b): the correlation between observed and predicted values was \nrLeWei{} for Wei et al. and \nrLeLiu{} for Liu et al., against \nrLeRo{} for the physics-residual model, and the mean bias was \nbiasLeAllWei{} and \nbiasLeAllLiu{} cm, against \nbiasLeAllRo{} cm. The bias of all series increased towards thin active layers (Fig.~\ref{fig:products}c): for site-years with an observed ALT below 30 cm (\nthinNLe{} in the Lena Delta), Wei et al. and Liu et al. overestimated by \nbiasLeThinWei{} and \nbiasLeThinLiu{} cm and the physics-residual model by \nbiasLeThinRo{} cm; in Alaska the same class (\nthinNAk{} site-years) was overestimated by \nbiasAkThinWei{}, \nbiasAkThinLiu{} and \nbiasAkThinRo{} cm.

\paperfig{Fig5_products}{1.0}{\textbf{Comparison with published products on the same cells and years.} \textbf{a}, RMSE of three methods and five products on the common site-years (\nyrscc), for all cells (filled) and for cells more than 25 km from the training sites of Wei et al. (open); values above 45 cm are drawn on the compressed segment. \textbf{b}, Observed against predicted ALT on the \nccLeN{} common site-years of the Lena Delta, with the 1:1 line and the Pearson correlation. \textbf{c}, Mean bias within classes of observed ALT (site-years per class below the axis; classes with fewer than 10 site-years omitted). \textbf{d}, RMSE per observation year for years with at least 30 site-years, with the number of site-years per year. Our methods use spatial-block out-of-fold predictions throughout.}{fig:products}

The regional contrast follows the training data of the products. Alaska label cells lie at a median distance of \ndistweiAkmed{} km from the training sites of Wei et al., and Liu et al. were trained on an earlier version of the same Alaska and Canada field data, whereas the Lena Delta cells lie at a median distance of \ndistweiLemed{} km (Supplementary Fig.~S6). Both products overestimate the thin active layers of the northern delta, with mean biases of \nbiasLeNorthWei{} and \nbiasLeNorthLiu{} cm north of 72$^\circ$N. A model trained without any Lena Delta label was already more accurate there than Wei et al. on the same site-years (\ntlvprmseLeRsWeizero{} versus \ntlvprmseLeWeizero{} cm; $\Delta$RMSE \ntlvpLeRsWeizero{} cm), so that the advantage in this region does not depend on local labels alone.

Maps of 2013, the year with the most observed cells in Alaska and Canada among the years observed in all three regions, show these differences in space (Fig.~\ref{fig:prodmaps}). At the \nfeAkN{} Alaska cells observed in that year, the RMSE of the three maps was similar (\nfeAkOof{}, \nfeAkWei{} and \nfeAkLiu{} cm), whereas the correlation with the observations was lower for our model (\nfeAkROof) than for the products (\nfeAkRWei{} and \nfeAkRLiu), \nfeAkNearWei{} of these cells lying within 25 km of a training site of Wei et al. The largest differences between the maps occur above 600 m in mountain terrain; between 68 and 70$^\circ$N the products are deeper than our map there by \nmapAkFootMtnWei{} and \nmapAkFootMtnLiu{} cm. In the Lena Delta both products were deeper than the observations (mean bias \nfeLeBiasWei{} and \nfeLeBiasLiu{} cm; RMSE \nfeLeWei{} and \nfeLeLiu{} cm against \nfeLeOof{} cm), although only \nfeLeN{} cells were observed in 2013. In Canada (NW), the RMSE at the \nfeCaN{} observed cells was \nfeCaOof{} cm for our model, \nfeCaWei{} cm for Wei et al. (mean bias \nfeCaBiasWei{} cm) and \nfeCaLiu{} cm for Liu et al., whose training data include an earlier version of these observations.

\paperfig{Fig6_product_maps}{1.0}{\textbf{ALT maps of 2013 from the physics-residual model and two published products, with the observed values.} \textbf{a--c}, Alaska: physics-residual model, Wei et al. and Liu et al.; the product maps are shown on the cells of our map, and the cells observed in 2013 are drawn as circles filled with the observed ALT on the same colour scale. \textbf{d}, Observed against predicted ALT at the \nfeAkN{} observed cells with values in all three maps, with the 1:1 line, the correlation and the RMSE; our predictions are spatial-block out-of-fold values. \textbf{e--h}, The same for the Lena Delta (\nfeLeN{} cells). \textbf{i--l}, The same for Canada (NW) (\nfeCaN{} cells). Each region has one colour range from the 5th--95th percentiles of the three maps inside the area of applicability; lighter colours mark cells outside it.}{fig:prodmaps}

Year by year, the physics-residual model had a lower error than both Wei et al. and the Stefan model in all \nywLen{} Lena Delta years with at least 30 site-years (\nyrsywLe), in \nywAkRoWeiwin{} of \nywAkRoWein{} Alaska years against Wei et al. and \nywAkRoPowin{} of \nywAkRoPon{} against the Stefan model (\nyrsywAk), and in \nywCaRoPowin{} of \nywCaRoPon{} Canadian years against the Stefan model (Fig.~\ref{fig:products}d). At Alaska cells observed in several years, the correlation between observed and predicted year-to-year changes was \nrdeltaAkRo{} for the physics-residual model and \nrdeltaAkWei{} and \nrdeltaAkLiu{} for Wei et al. and Liu et al., whereas in the Lena Delta the products followed the interannual changes more closely (\nrdeltaLeWei{} and \nrdeltaLeLiu{} versus \nrdeltaLeRo).

The residual learner can also be applied with a product as the anchor (equation~(\ref{eq:prod})). With the Lena Delta labels, adding the residual learner to the recalibrated product of Wei et al. changed its error by \npadLeRpPcal{} cm, and the Stefan-anchored model remained more accurate than the product-anchored model (\npadLeRoRp{} cm; Supplementary Fig.~S7); in Alaska and Canada the two anchors were not distinguishable.

\subsection{Year-labelled 1 km ALT maps, 2011--2024}
The physics-residual models fitted to all site-years of each region were applied with the ERA5-Land climate of each year from 2011 to 2024 (Fig.~\ref{fig:maps}); the maps cover \nncellmapAk{} cells in Alaska, \nncellmapLe{} in the Lena Delta and \nncellmapCa{} in Canada (NW). Relative to the 2011--2024 mean, the regional mean anomaly ranged from \nanminAk{} cm in \nanminyrAk{} to \nanmaxAk{} cm in \nanmaxyrAk{} in Alaska, from \nanminLe{} cm in \nanminyrLe{} to \nanmaxLe{} cm in \nanmaxyrLe{} in the Lena Delta and from \nanminCa{} cm in \nanminyrCa{} to \nanmaxCa{} cm in \nanmaxyrCa{} in Canada, with spatial patterns that differ between years (Fig.~\ref{fig:maps}f). In the Lena Delta, the mean anomaly of the observed cells had the same sign as the predicted regional anomaly in \nanobssameLe{} of \nanobsnLe{} years with at least 50 site-years (\nanobsyrsLe; $r$ = \nanobsrLe), and the correlation remained \nanobsrLeNoTwenty{} without 2020, when the observed anomaly of \nanobsLeTwentyN{} site-years reached \nanobsLeTwentyObs{} cm. In Alaska, where the set of observed cells changes between campaigns, the agreement was weaker ($r$ = \nanobsrAk) and the predicted interannual range (\nanpredrangeAk{} cm) was smaller than that of the observed cell means (\nanobsrangeAk{} cm); in Canada the correlation over \nanobsnCa{} years was \nanobsrCa, with the observed anomalies of 2023 and 2024 (80 site-years each) far above the predicted range (Fig.~\ref{fig:maps}d). The linear trend of the yearly maps averaged \ntrendAk{}, \ntrendLe{} and \ntrendCa{} cm per decade over the Alaska, Lena Delta and Canada cells, with positive values in \ntrendposAk\%, \ntrendposLe\% and \ntrendposCa\% of the cells (Fig.~\ref{fig:maps}e); it reflects the ERA5-Land climate of the period as transmitted by the model and is not an observed trend. The area of applicability covers \naoaAk\% of the Alaska cells, \naoaLe\% of the Lena Delta cells and \naoaCa\% of the Canada cells; cells outside it lie mainly in mountains and at the southern margin of the permafrost zone (Supplementary Fig.~S8). The half-width of the 90\% prediction interval is \nhwAk{}, \nhwLe{} and \nhwCa{} cm. The correction added by the residual learner is mostly within $\pm$10 cm (Supplementary Fig.~S9), and series at long-record sites are shown in Supplementary Fig.~S10.

\paperfig{Fig7_maps}{1.0}{\textbf{Year-labelled 1 km ALT maps, 2011--2024.} \textbf{a--c}, Mean ALT of 2011--2024 from the yearly maps of Alaska, the Lena Delta and Canada (NW), with windows enlarged at the right (1 km cells) and the observed cells filled with their observed mean ALT; lighter colours mark cells outside the area of applicability. \textbf{d}, Regional mean anomaly of the maps and mean anomaly of the observed cells in years with at least 50 site-years, centred on the predicted anomaly of the same years (y ranges differ between regions). \textbf{e}, Linear trend of the yearly maps (least squares per cell), which describes the modelled response to the climate of each year. \textbf{f}, ALT anomaly of each year relative to the 2011--2024 mean in Alaska.}{fig:maps}
